\subsection{TRACE OF AN ENDOMORPHISM}

Let C be a commutative ring and E a C-module. The mapping \((x^{*}, x) \mapsto \langle x, x^{*} \rangle\) of \(E^{*} \times E\) into C is then C-bilinear, since, for all \(\gamma \in C\) , \(\langle \gamma x, x^{*} \rangle = \gamma \langle x, x^{*} \rangle\) and \(\langle x, x^{*} \gamma \rangle = \langle x, x^{*} \rangle \gamma\) ; we derive a canonical C-linear mapping

\[
\tau : \mathbf {E} ^ {*} \otimes_ {\mathbf {c}} \mathbf {E} \rightarrow \mathbf {C}\tag{16}
\]

such that \(\tau(x^{*} \otimes x) = \langle x, x^{*} \rangle\) (§ 3, no. 5). Suppose now also that E is a finitely generated projective C-module; the canonical isomorphism (11) of no. 2 is then a C-module isomorphism and we can therefore define by transporting the structure a canonical linear form \(\mathrm{Tr} = \tau \circ \theta_{\mathbf{E}}^{-1}\) on the C-module \(\operatorname{End}_{\mathbf{C}}(\mathbf{E})\) . For all \(u \in \operatorname{End}_{\mathbf{C}}(\mathbf{E})\) the scalar \(\operatorname{Tr}(u)\) is called the trace of the endomorphism \(u\) ; every \(u \in \operatorname{End}_{\mathbf{C}}(\mathbf{E})\) can be written (in general in an infinity of ways) in the form \(x \mapsto \sum_{i} \langle x, x_{i}^{*} \rangle y_{i}\) where \(x_{i}^{*} \in \mathbf{E}^{*}\) and \(y_{i} \in \mathbf{E}_{i}\) by virtue of no. 2, Corollary to Proposition 2; then

\[
\operatorname{Tr} (u) = \sum_ {i} \left\langle y _ {i}, x _ {i} ^ {*} \right\rangle \quad (\text { cf. } \S 1 0, \text { no. } 1 1).\tag{17}
\]

By definition,

\begin{enumerate}
\def\labelenumi{(\arabic{enumi})}
\setcounter{enumi}{17}
\tightlist
\item
\end{enumerate}

\[
\operatorname{Tr} (u + v) = \operatorname{Tr} (u) + \operatorname{Tr} (v)\tag{19}
\]

\[
\operatorname{Tr} (\gamma u) = \gamma \operatorname{Tr} (u)
\]

for \(u, v\) in \(\operatorname{End}_{\mathbf{C}}(\mathbf{E})\) and \(\gamma \in \mathbf{C}\) . Moreover:

PROPOSITION 3. Let C be a commutative ring, E, F two finitely generated projective C-modules and \(u: \mathbf{E} \to \mathbf{F}\) and \(v: \mathbf{F} \to \mathbf{E}\) two linear mappings; then

\[
\operatorname{Tr} (v \circ u) = \operatorname{Tr} (u \circ v).\tag{20}
\]

The two mappings \((u, v) \mapsto \operatorname{Tr}(u \circ v)\) , \((u, v) \mapsto \operatorname{Tr}(v \circ u)\) of

\[
\operatorname{Hom} _ {\mathbf {c}} (\mathbf {E}, \mathbf {F}) \times \operatorname{Hom} _ {\mathbf {c}} (\mathbf {F}, \mathbf {E})
\]

into C are C-bilinear; it therefore suffices to verify (20) when u is of the form \(x \mapsto \langle x, a^{*} \rangle b\) and v of the form \(y \mapsto \langle y, b^{*} \rangle a\) , with \(a \in E\) , \(a^{*} \in E^{*}\) , \(b \in F\) , \(b^{*} \in F^{*}\) . But then \(v \circ u\) is the mapping \(x \mapsto \langle x, a^{*} \rangle \langle b, b^{*} \rangle a\) and \(u \circ v\) the mapping \(y \mapsto \langle y, b^{*} \rangle \langle a, a^{*} \rangle b\) . Formula (17) shows that the values of the two sides of (20) are equal to \(\langle a, a^{*} \rangle \langle b, b^{*} \rangle\) .

COROLLARY. If \(u_1, \ldots, u_p\) are endomorphisms of \(\mathbf{E}\) , then

\[
\operatorname{Tr} \left(u _ {1} \circ u _ {2} \circ \dots \circ u _ {p}\right) = \operatorname{Tr} \left(u _ {i} \circ u _ {i + 1} \circ \dots \circ u _ {p} \circ u _ {1} \circ \dots \circ u _ {i - 1}\right)
\]

for \(1 \leqslant i \leqslant p\) (``invariance of the trace under cyclic permutation'').

It suffices to apply (20) to the product

\[
\left(u _ {1} \circ u _ {2} \circ \dots \circ u _ {i - 1}\right) \circ \left(u _ {i} \circ u _ {i + 1} \circ \dots \circ u _ {p}\right).
\]

Note that on the other hand it is not necessarily true that

\[
\operatorname{Tr} (u \circ v \circ w) = \operatorname{Tr} (u \circ w \circ v)
\]

for three endomorphisms u, v, w of E.

